% Four-point unit-weight example, with reference orientations as in chapter 14.
\begin{tikzpicture}[x=1mm,y=1mm,>=Stealth,
  font=\figurefont\fontsize{8}{10}\selectfont,text=FigureInk,
  vertex/.style={circle,draw=FigureStroke,fill=FigureBlueFill,
    line width=.65pt,minimum size=6mm,inner sep=0pt},
  labelbox/.style={inner sep=1pt,fill=white},
  ordinary/.style={draw=FigureStroke,line width=.65pt},
  directed/.style={ordinary,->},
  transfer/.style={draw=FigureBlue,line width=.8pt,->}]
  \def\leftpanel{a}\def\middlepanel{b}\def\rightpanel{c}
  \foreach \offset/\panel in {0/a,37/b,74/c}{
    \begin{scope}[xshift=\offset mm]
    \node[vertex] (v1) at (0,18) {$1$};
    \node[vertex] (v2) at (18,18) {$2$};
    \node[vertex] (v3) at (18,0) {$3$};
    \node[vertex] (v4) at (0,0) {$4$};
    \ifx\panel\middlepanel
      \draw[directed] (v1)--node[above,labelbox] {$1$} (v2);
      \draw[directed] (v2)--node[right,labelbox] {$1$} (v3);
      \draw[directed] (v3)--node[below,labelbox] {$1$} (v4);
      \draw[directed] (v4)--node[left,labelbox] {$-3$} (v1);
      \draw[directed] (v1)--node[above right,labelbox,inner sep=.3pt] {$2$} (v3);
    \else
      \draw[ordinary] (v1)--(v2)--(v3)--(v4)--(v1) (v1)--(v3);
    \fi
    \ifx\panel\leftpanel
      \node[left=1mm of v1,labelbox] {$3$};
      \node[right=1mm of v2,labelbox] {$2$};
      \node[right=1mm of v3,labelbox] {$1$};
      \node[left=1mm of v4,labelbox] {$0$};
    \fi
    \ifx\panel\rightpanel
      \node[left=1mm of v1,labelbox] {$6$};
      \node[right=1mm of v2,labelbox] {$0$};
      \node[right=1mm of v3,labelbox] {$-2$};
      \node[left=1mm of v4,labelbox] {$-4$};
    \fi
    \node at (9,-10) {(\panel)};
    \end{scope}
  }
  \draw[transfer] (23,9)--node[above] {$\bm B^{\top}$} (32,9);
  \draw[transfer] (60,9)--node[above] {$\bm B$} (69,9);
\end{tikzpicture}
